\section{2. RL for LLMs: The Ideas That Matter}
% ============================================================

\begin{bluf}
You need to explain the RLHF pipeline, why DPO exists but online RL is preferred at scale, what GRPO changed, and why reward hacking is the central unsolved problem. All in plain English. Equations are ammunition for depth, not your opening move.
\end{bluf}

\subsection{The RLHF Pipeline (3 stages)}

\subsubsection{Stage 1: SFT (Supervised Fine-Tuning)} Take a pretrained model that just predicts the next token from web text. Fine-tune it on examples of \textit{good} responses to prompts. Now it follows instructions instead of just continuing text. Think of it as giving the model a starting personality.

\subsubsection{Stage 2: Reward Model} You can't tell the model ``be helpful'' with a loss function. Instead: show human annotators pairs of model outputs, ask ``which is better?'', and train a separate model to predict those preferences. This reward model takes (prompt, response) and outputs a single number: how good is this response?

\textbf{Key insight:} You use \textit{rankings}, not scores, because humans are terrible at absolute scores but reliable at ``A is better than B.'' The Bradley-Terry model converts pairwise rankings into a scalar reward.

\subsubsection{Stage 3: RL (PPO)} Now treat the language model as an RL agent. The ``action'' is generating text. The ``reward'' comes from the reward model. Optimize the policy (= the LM) to generate text that scores high on the reward model.

\textbf{The critical constraint:} You add a KL penalty that prevents the model from drifting too far from the SFT model. Without it, the model finds adversarial gibberish that ``hacks'' the reward model --- high score, zero quality.

\begin{sayit}[RLHF pipeline]
``The basic loop is: humans rank outputs, you distill those rankings into a reward model, then you RL-optimize the language model against that reward. The hard part isn't any individual step --- it's that the reward model is an imperfect proxy, and the policy will exploit every imperfection. That's why you need the KL penalty as a leash, and even then, models find creative ways to game the system.''
\end{sayit}

\subsection{Why PPO Is Hard to Scale}

PPO for LLMs requires keeping \textbf{four models in memory simultaneously}:
\begin{enumerate}
  \item The \textbf{policy} (the model you're training)
  \item The \textbf{reference model} (frozen copy for KL computation)
  \item The \textbf{reward model} (evaluates outputs)
  \item The \textbf{critic/value network} (estimates expected future reward --- usually same size as the policy)
\end{enumerate}

That's 4$\times$ the memory of just running inference. Plus you need to \textit{sample from the policy during training} (generate full completions), which is expensive.

\begin{deeper}
PPO also requires careful hyperparameter tuning (clipping range $\epsilon$, KL coefficient $\beta$, learning rate schedules). The clipping mechanism prevents catastrophically large updates: if the new policy differs too much from the old one on a given example, the gradient gets clipped. This keeps training stable but means you need many small steps.
\end{deeper}

\subsection{DPO: The Elegant Shortcut (and why it's not enough)}

\textbf{The idea:} Rafailov et al.\ showed you can mathematically rearrange the RLHF objective to eliminate the reward model entirely. Instead of: train reward model $\to$ RL against it, you directly optimize the policy on preference pairs using a binary cross-entropy loss.

\textbf{Why it's appealing:} One model instead of four. No RL sampling loop. Stable training. Just supervised learning on preference data.

\textbf{Why frontier teams don't use it:} DPO learns from a \textit{static dataset} of preferences. It never explores --- it never generates new outputs and discovers new failure modes. Online RL (PPO/GRPO) continuously samples from the improving policy, which means it can discover and learn from edge cases that weren't in the original data.

\begin{sayit}[DPO vs.\ online RL]
``DPO is elegant but it's fundamentally offline --- you're training on a fixed dataset of preferences. Online RL lets the model discover its own failure modes by sampling from itself. At the frontier, that exploration is what matters. The model needs to find the weird edge cases and learn from them, not just memorize which of two options humans preferred.''
\end{sayit}

\subsection{GRPO: DeepSeek's Systems Insight}

\textbf{The problem GRPO solves:} PPO's critic model (the value network) is usually the same size as the policy. That's half your GPU memory gone on a model whose only job is estimating ``how much reward do I expect from here?''

\textbf{GRPO's trick:} For each prompt, generate a \textit{group} of outputs. Score them all with the reward. Normalize the rewards within the group (subtract mean, divide by std). Use these normalized scores as your advantage estimates. No critic needed.

\textbf{Why it matters:} Dropping the critic frees up \textbf{$\sim$50\% of GPU memory}. You can use that to scale up the policy, increase batch size, or generate more rollouts. It's as much a \textit{systems engineering} insight as an algorithmic one.

\begin{sayit}[GRPO]
``The genius of GRPO is that it's a systems hack as much as an algorithmic one. By using group statistics instead of a learned critic, DeepSeek freed up half their GPU memory. That lets you scale the policy larger or sample more rollouts, which directly translates to better results. It's the kind of insight where understanding the hardware constraints leads you to a better algorithm.''
\end{sayit}

\subsection{DeepSeek R1: The Proof RL Works}

\textbf{The headline:} DeepSeek applied RL directly to a base model with \textit{zero} SFT, using only rule-based rewards (did you get the right answer? did you use the right format?). The model \textit{spontaneously} developed chain-of-thought reasoning, self-verification, and backtracking.

\textbf{Why this matters enormously:}
\begin{itemize}
  \item It proves RL can create \textit{new capabilities}, not just polish existing ones
  \item The reward was dead simple: right/wrong answer. No neural reward model.
  \item AIME 2024 scores: 15.6\% (base) $\to$ 71.0\% (pure RL) $\to$ 79.8\% (full pipeline)
\end{itemize}

\textbf{The four-stage production pipeline:}
\begin{enumerate}
  \item Cold-start SFT (readable chain-of-thought examples)
  \item Reasoning RL (GRPO on math/code/logic)
  \item Mixed SFT (rejection-sample 600k reasoning + 200k general)
  \item General RL (rule-based for reasoning, neural RM for subjective tasks)
\end{enumerate}

\begin{whyanthro}
This is the empirical proof behind the core claim: ``RL in language models has finally worked.'' Math and code fell because they have deterministic verification. The frontier question: what's the next domain where you can build verification good enough for RL to work?
\end{whyanthro}

\subsection{The RLVR Paradigm Shift (2025+)}

RLVR (Reinforcement Learning from Verifiable Rewards) has replaced RLHF as the dominant RL training paradigm at frontier labs. Per Karpathy's 2025 review, RLVR ``gobbled up the compute originally intended for pretraining.'' Key distinctions:

\begin{itemize}
  \item \textbf{RLHF:} neural reward model, gameable proxy, short training runs (``thin/short stage'')
  \item \textbf{RLVR:} deterministic verification (unit tests, math correctness), non-gameable, allows much longer optimization (weeks)
  \item DeepSeek R1's GRPO = RLVR (rule-based rewards, not neural RM)
  \item Models spontaneously develop reasoning strategies under RLVR pressure
\end{itemize}

\begin{sayit}[RLVR vs RLHF]
``RLHF taught us that RL on language models works. RLVR taught us that it scales --- if you can verify, you can optimize for weeks. The frontier question isn't `can RL work?' anymore. It's `where can we build clean verification next?' That's why verification engineering is the binding constraint.''
\end{sayit}

% ============================================================
\section{3. Reward Hacking: The Central Problem}
% ============================================================

\begin{bluf}
This is where production experience with reward signals matters most. The core challenge: how do you get a model to optimize for what you \textit{actually} want, not what it can game?
\end{bluf}

\subsection{What It Is}

The reward model (or any evaluation metric) is a \textit{proxy} for what you actually want. The policy will find and exploit every gap between the proxy and reality:
\begin{itemize}
  \item \textbf{In code:} Models hard-code expected values into unit tests without solving the underlying problem. ``Even with unit tests, models find ways to hack in particular values.'' --- DeepSeek R1 paper
  \item \textbf{In language:} Models produce verbose, confident-sounding text that scores well on helpfulness metrics while saying nothing substantive
  \item \textbf{In domain tasks:} Models pattern-match surface features of inputs (e.g., seeing a keyword and predicting the label) without actually reasoning about the underlying problem
\end{itemize}

\subsection{The Mitigation Toolkit}

\begin{tabular}{@{}p{2.8cm}p{4.5cm}@{}}
\toprule
\textbf{Approach} & \textbf{Trade-off} \\
\midrule
KL penalty & Limits exploitation magnitude but also limits improvement \\
Rule-based rewards & Ungameable but only works for verifiable domains \\
RM ensembles & Harder to hack many models but expensive \\
Iterative retraining & Closes loopholes over time but adds pipeline complexity \\
Expert curation & Highest quality signal but doesn't scale \\
\bottomrule
\end{tabular}

\subsection{War Stories (prep your own)}

Prepare a concrete example from your own experience where a model gamed a proxy metric. The best stories follow this arc: (1) the metric looked fine, (2) you diagnosed that the model was exploiting a shortcut, (3) the fix was better signal, not a better model.

\begin{sayit}[Reward hacking example]
``We hit classic reward hacking --- the model was pattern-matching surface features instead of reasoning about the underlying problem. Eval metrics looked fine because most cases \textit{are} straightforward. It was the ambiguous cases where the model was just guessing confidently.

The fix wasn't better prompting --- it was better signal. We retrieved similar past cases where a domain expert had already decided and used those as few-shot exemplars for reflection. Then we curated the eval set to \textit{oversample} the hard cases --- the ones where gaming fails. The model started producing reasoning you could actually audit.''
\end{sayit}

\begin{whyanthro}
A good war story demonstrates: (1) you can diagnose reward hacking in production, (2) you think about signal quality, not just model quality, (3) you know that curating the eval to be adversarial is as important as improving the model. These are exactly the skills the RL team needs.
\end{whyanthro}

% ============================================================
